\documentclass[11pt]{article}
\usepackage[margin=2.0cm]{geometry}
\usepackage{graphicx,amsmath,booktabs,hyperref,float,caption}
\hypersetup{colorlinks=true,linkcolor=blue,urlcolor=blue}
\captionsetup{font=small}
\newcommand{\crad}{C^{\rm rad}}
\newcommand{\Lum}{\mathcal{L}}
\newcommand{\BR}{\mathcal{B}}

\title{Extracting the Born cross section for $ep\to e'p'J/\psi$ at CLAS12:\\
procedure, master formula, and the traps in our own samples\\[4pt]
\large Campaign \texttt{prod\_jpsi\_f18in\_20261009}}
\author{\texttt{exclurad\_py} note, branch \texttt{jpsi-born-xsec-note}}
\date{6 November 2026}

\begin{document}
\maketitle

\begin{abstract}
What to do with a data yield binned in $(W,Q^2,t)$, a radiative Monte Carlo and a Born Monte Carlo
in order to obtain a Born cross section. The master formula is given in two algebraically equivalent
forms. The one to \emph{compute} is the ratio of the data yield to the reconstructed radiative
Monte Carlo, times the model: in that form a wrong model \emph{normalisation} cancels exactly, which
matters here because \texttt{jpsi.diffrad\_vm} sits a factor 4--6 below the data. Only the shape
inside the bin survives, and it is iterated away. Six inputs beyond the three samples are required
and are listed explicitly; the branching ratio $\BR(J/\psi\to e^+e^-)$ is one of them and is
\emph{not} applied anywhere in the generator, a factor 16.7 if missed. Five properties of the
existing campaign were verified while writing this note and are reported in Sec.~5, including the
fact that the radiative and Born cards differ in three physics fields rather than one, that the Born
outbending arm is 28\,\% short of the radiative arm, and that the generator's \texttt{sigma\_nb}
cannot be used to form the radiative correction.
\end{abstract}

\section{The inputs}

Three samples are not enough. The complete set is Table~\ref{tab:inputs}. Items 4--10 are not
refinements: each is a factor, or a bias of the order of the effect being measured.

\begin{table}[H]
\centering\small
\begin{tabular}{clp{7.6cm}}
\toprule
\# & Input & Source \\
\midrule
1 & data yield in $(W,Q^2,t)$ bins & analysis ntuple \\
2 & radiative MC through GEMC & \texttt{prod\_jpsi\_f18in\_20261009/rad} \\
3 & Born MC through GEMC & \texttt{prod\_jpsi\_f18in\_20261009/born} \\
4 & integrated luminosity $\Lum$ & for exactly the run/trigger selection that made the yield \\
5 & signal extraction prescription & $M(ee)$ window and background shape \\
6 & $\BR(J/\psi\to e^+e^-)=(5.971\pm0.032)\,\%$ & PDG. \textbf{Not in the generator}, Sec.~5.1 \\
7 & data/MC efficiency scale factors & FT electron, lepton and proton tracking, trigger \\
8 & bin volumes, centroids, averaging rule & chosen once and written down \\
9 & one cut list for data and \emph{both} MC arms & exclusivity cut included \\
10 & a model to iterate with & Sec.~3, step 7 \\
\bottomrule
\end{tabular}
\caption{Inputs to the extraction.}
\label{tab:inputs}
\end{table}

\section{The master formula}

Each arm carries a per-event weight built from what the generator already reports in
\texttt{<job>.config.json},
\begin{equation}
w^{\rm arm}=\frac{\sigma_{\rm gen}^{\rm arm}}{N_{\rm gen}^{\rm arm}}\quad[\mathrm{nb/event}],
\qquad \sigma_{\rm gen}\equiv\texttt{report.sigma\_nb}.
\end{equation}
Let $n_i$ be the number of \emph{reconstructed} Monte Carlo events of a given arm falling in bin $i$
after the identical analysis chain, and $m_i$ the number of \emph{generated} (truth) events in bin
$i$ with no cuts at all.

\paragraph{Form A --- ratio to the radiative MC. This is the one to compute.}
\begin{equation}
\boxed{\;\sigma_i^{\rm Born}=R_i\,\bar\sigma_i^{\rm Born,model},
\qquad
R_i=\frac{N_i^{\rm data}}{\Lum\,\BR\,w^{\rm rad}\,n_i^{\rm rad}}\;}
\label{eq:formA}
\end{equation}

\paragraph{Form B --- explicit factors. This is the one to tabulate systematics on.}
\begin{equation}
\bar\sigma_i^{\rm Born}
=\frac{N_i^{\rm data}}{\Lum\,\BR\,\Delta\Omega_i\,\varepsilon_i^{\rm Born}\,\crad_i},
\qquad
\varepsilon_i^{\rm Born}=\frac{n_i^{\rm born}}{m_i^{\rm born}},
\qquad
\crad_i=\frac{w^{\rm rad}\,n_i^{\rm rad}}{w^{\rm born}\,n_i^{\rm born}}.
\label{eq:formB}
\end{equation}

Equations~(\ref{eq:formA}) and~(\ref{eq:formB}) are the same once the bin-averaging is done
consistently. Form~A is what should be computed, for three reasons.

\begin{enumerate}
\item Acceptance, efficiency, resolution smearing, bin migration and radiation all sit inside one
measured number, $n_i^{\rm rad}$, obtained by running the same analysis code over reconstructed
Monte Carlo. Nothing is factorised that does not factorise.

\item The $M(ee)$ window efficiency cancels. The same window is applied to the data and to
$n_i^{\rm rad}$, so it never has to be computed, and the radiative lineshape is simulated rather
than assumed.

\item \textbf{A wrong model normalisation cancels exactly.} If the model is scaled by a factor $f$,
then $n_i^{\rm rad}\to f\,n_i^{\rm rad}$, hence $R_i\to R_i/f$ and
$\bar\sigma_i^{\rm Born,model}\to f\,\bar\sigma_i^{\rm Born,model}$, so
\begin{equation}
\sigma_i^{\rm Born}=R_i\,\bar\sigma_i^{\rm Born,model}
\quad\text{is invariant under}\quad \sigma^{\rm model}\to f\,\sigma^{\rm model}.
\end{equation}
\end{enumerate}

The third point is not academic. The model \texttt{jpsi.diffrad\_vm} carries
$\sigma_T\propto Q^{-3.11}$, which double-counts the $1/Q^2$ already present in the virtual-photon
flux, and it lies a factor 4--6 below the data at Forward Tagger kinematics. Form~A is immune to
that. What survives is only the \emph{shape inside the bin}, and step~7 iterates it away.

\section{Procedure}

\paragraph{Step 1 --- one cut list.}
Write the cuts once and apply them to the data, the radiative MC and the Born MC from the same code.
Include the exclusivity cut. Both cards carry \texttt{vcut = 0}, so nothing was cut at generation
and the whole radiative phase space is in the files; the exclusivity cut is applied \emph{offline},
and that is the only way migration through resolution is simulated.

\paragraph{Step 2 --- signal yield per bin.}
Count in a window around $M(ee)=3.0969$~GeV. Do \emph{not} fit a free lineshape: a free Gaussian
width absorbs the radiative tail and biases the yield. Window counting reduces that model bias by
about a factor 5. Subtract background from sidebands or a fitted continuum, and propagate the
subtraction error into $N_i^{\rm data}$.

\paragraph{Step 3 --- the same window on the MC.}
The identical $M(ee)$ window applied to the reconstructed radiative MC gives $n_i^{\rm rad}$. The
radiative arm runs \texttt{idecay = 2}, so the decay final-state radiation is simulated and the
lineshape is never assumed; that is what makes the window efficiency cancel between the data and
$n_i^{\rm rad}$. The Born arm runs \texttt{idecay = 1}, no decay radiation at all, which is why
$\crad$ built from these twins also carries the decay-lepton correction (Sec.~5.2).

\paragraph{Step 4 --- normalise.}
$Y_i^{\rm rad}=\Lum\,\BR\,w^{\rm rad}\,n_i^{\rm rad}$. Average \texttt{sigma\_nb} over jobs; all
production jobs share the card, so $w^{\rm rad}$ is common to them.

\paragraph{Step 5 --- form $R_i$ and multiply by the model.}
Bin-average the model's Born cross section over the bin using the model's own density, \emph{or}
quote at the bin centroid with a bin-centring correction. State which. The two differ wherever the
$t$-slope is steep, which near threshold is everywhere.

\paragraph{Step 6 --- compute $\crad_i$ in any case.}
Form~A does not need it, but it must be reported, and its bin-to-bin behaviour is the diagnostic
that says whether the correction is sane. On this Monte Carlo we measure $0.80/0.76/0.62$ by $W$
bin, falling as $W$ rises. A $\crad$ that is flat in $W$, or above unity, means something upstream
is wrong.

\paragraph{Step 7 --- iterate.}
Reweight the Monte Carlo event by event by $\sigma^{\rm extracted}/\sigma^{\rm model}$ as a function
of $(W,Q^2,t)$, then recompute $n_i^{\rm rad}$ and $R_i$. Two or three passes converge. This is
required rather than cosmetic: the model is fitted over $Q^2=0.015$--$0.1$~GeV$^2$ while the
generation box runs $0.001$--$0.25$, so without iteration the bin-averaging is model-dominated at
both $Q^2$ edges.

\paragraph{Step 8 --- efficiency scale factors.}
GEMC does not reproduce the Forward Tagger electron or the tracking efficiencies. Apply the measured
data/MC ratios to $Y_i^{\rm rad}$, per bin wherever they are kinematics-dependent.

\paragraph{Step 9 --- quote it.}
For $ep\to e'p'J/\psi$ no flux factor is needed. If converting to
$\sigma(\gamma^*p\to J/\psi\,p)$, state the flux convention (Hand or Gilman) explicitly: it moves
the number.

\section{Statistical error}

From Form~A, with the three counts independent,
\begin{equation}
\left(\frac{\delta R_i}{R_i}\right)^{2}
=\frac{1}{N_i^{\rm data}}+\frac{1}{n_i^{\rm rad}}
+\left(\text{background subtraction term}\right).
\end{equation}
The Born arm enters the statistical error \emph{only} through $\crad$ and
$\varepsilon^{\rm Born}$, that is only in Form~B and in the reported diagnostics. That is not a
reason to generate less Born: see Sec.~5.4.

\section{Facts checked in the actual samples, 6 November 2026}

\subsection{The branching ratio is not in the generator}

\begin{verbatim}
grep -rniE "branching|0.0597|5.97" --include=*.py exclurad_py/   ->   nothing
\end{verbatim}
\texttt{report.sigma\_nb} is the \textbf{full $J/\psi$ production cross section}. Every event is
decayed to $e^+e^-$ with probability one (\texttt{idecay = 2}). The data yield is an $e^+e^-$ yield.
Therefore $\BR$ multiplies the Monte Carlo, once. Omitting it makes the answer wrong by 16.7.

\subsection{The twins differ in three physics fields, not one}

From \texttt{prod\_jpsi\_f18in\_20261009/\{born,rad\}/prov/*\_043.card.txt}:
\begin{center}
\begin{tabular}{llccl}
\toprule
field & meaning & born & rad & switches \\
\midrule
4 & \texttt{iborn} & 1 & 0 & vertex radiation \\
12 & \texttt{idecay} & 1 & 2 & decay-lepton radiation \\
18 & \texttt{itarg} & 0 & 1 & pre-vertex radiator \\
\bottomrule
\end{tabular}
\end{center}
plus the seed (field 11) and \texttt{outroot} (field 13), which differ as they must.

All three are \emph{correct}: Born means no radiation of any kind, so $\crad$ built from these twins
is the \emph{whole} correction --- electron-line radiation, straggling in the target before the
vertex, and radiation of the decay leptons. This is also what
\texttt{docs/rc\_method\_jpsi/rc\_method\_jpsi\_report.tex} states in its abstract; the first
revision of the present note said ``two fields'' and had missed \texttt{idecay}.

The card header, however, claims ``IDENTICAL except iborn (field 4)'' and offers a \texttt{diff}
command to prove it. That command in fact prints five differences. The danger is not the confusion
but the repair: a reader who trusts the comment may ``fix'' field~12 or field~18 to make the diff
clean, and a Born arm with \texttt{itarg = 1} or \texttt{idecay = 2} silently reduces $\crad$ to the
vertex term alone. It is the comment that is wrong, not the cards.

The \texttt{itarg = 0} point-vertex bug is \textbf{not} present in this sample: the fix
\texttt{bbe6628} (2026-10-22 11:25) is an ancestor of the sample's recorded commit \texttt{8ac2fc0}
(2026-10-22 12:36), so the Born arm has the proper 5~cm vertex spread rather than $v_z=0$.

\subsection{$\crad$ cannot be taken from \texttt{sigma\_nb}}

\begin{center}
\begin{tabular}{lr}
\toprule
arm & \texttt{report.sigma\_nb} \\
\midrule
born & $3.494470\times10^{-4}$ \\
rad  & $3.368744\times10^{-4}$ \\
\midrule
ratio rad/born & $0.9640$ \\
\bottomrule
\end{tabular}
\end{center}
\texttt{sigma\_nb} counts events \emph{before} the acceptance preset, so any rad/Born ratio built
from it is about 1.33\,\% high, and in any case it is a box average rather than a bin correction.
$\crad_i$ must come from reconstructed yields inside the bin.

\subsection{The arms are not yet at equal statistics}

\begin{center}
\begin{tabular}{lrr}
\toprule
 & born & rad \\
\midrule
\texttt{out\_inb}  & 1120 & 1167 \\
\texttt{out\_outb} &  852 & 1180 \\
\bottomrule
\end{tabular}
\end{center}
The Born outbending arm is 28\,\% short, which caps the precision of both $\crad$ and
$\varepsilon^{\rm Born}$. It should be finished before a systematic is quoted on either.

\subsection{\texttt{iacc = jpsi} is a hard generation cut, not a flag}

\texttt{events.py} executes \texttt{continue} on failure: the event is never written. The windows
travel with the sample in the \texttt{.lund.info} sidecar:
\begin{center}
\begin{tabular}{lcc}
\toprule
particle & $|p|$ (GeV) & $\theta$ (deg) \\
\midrule
proton   & 0.2--4.5 & 3--45 \\
daughter & 1.0--9.0 & 5--40 \\
electron & 0.1--2.5 & 2--20 \\
\bottomrule
\end{tabular}
\end{center}
If any analysis cut is \emph{looser} than the corresponding preset edge, the Monte Carlo has exactly
zero events there and the acceptance is wrong, not merely imprecise. Measured in a 2000-event
sample, the live edges are the proton $\theta$ floor (minimum $3.013^\circ$, zero events below) and
the daughter $\theta$ ceiling (maximum $39.96^\circ$). The proton $|p|$ window never bites: the
kinematic minimum is $0.7156$~GeV/$c$ against a $0.2$~GeV/$c$ floor.

\subsection{$t$ from the proton, and the choice of $t$ or $t'$}

The proton $t$ estimator has RMS $0.047$~GeV$^2$ against $0.230$ for the $J/\psi$ one, a factor
5--6 better. The existing analysis ntuple carries \texttt{t\_prime}, not \texttt{t}. Since
$|t|_{\min}$ runs from $2.10$~GeV$^2$ at threshold to $0.4536$ at the top of the $W$ range, bins in
$t'$ and in $t$ map onto appreciably different regions. Fix the choice once and use it in the data,
in both Monte Carlo arms, and in the model bin-averaging.

For reference, the recoil proton follows from $t$ alone through $t=(p-p')^2=2M^2-2ME_p$, so
\begin{equation}
E_p=M+\frac{|t|}{2M},
\qquad
|\mathbf{p}_p|=\sqrt{|t|}\,\sqrt{1+\frac{|t|}{4M^2}} .
\label{eq:pproton}
\end{equation}
At the kinematic ceiling $W^2=M^2+2ME_{\rm beam}=20.7717$~GeV$^2$ ($W=4.5576$~GeV),
\begin{equation}
|t|_{\min}=0.4536~\mathrm{GeV}^2,
\qquad
|\mathbf{p}_p|=0.7156~\mathrm{GeV}/c,
\qquad
\theta_p=0.67^\circ ,
\end{equation}
so at the forward edge the recoil proton travels down the beamline and is undetectable; the
$3^\circ$ preset floor removes it. Over the whole box the recoil proton never exceeds
$\theta_p=25.3^\circ$: it is a Forward Detector particle everywhere, never a Central Detector one.

\section{Systematics}

\begin{table}[H]
\centering\small
\begin{tabular}{lp{9.5cm}}
\toprule
Source & How to estimate \\
\midrule
model shape inside the bin & spread over iterations; vary the $Q^2$ slope within its fit error \\
exclusivity cut & vary the $MX^2$ window; $\crad$ is the quantity that moves \\
signal extraction & window width; sideband versus fitted background \\
MC statistics & $n^{\rm rad}$ and $n^{\rm born}$ separately \\
efficiency scale factors & their own uncertainties, bin dependent \\
luminosity & from the data side \\
branching ratio & 0.54\,\% relative \\
bin centring & centroid versus bin-averaged model \\
\bottomrule
\end{tabular}
\caption{Systematic uncertainties to carry through Eq.~(\ref{eq:formB}).}
\end{table}

\section{Ntuple branches the chain needs}

The reader in \texttt{production/analysis/campaign\_plots.py} uses
\begin{verbatim}
E_photon, Q2, W, t_prime, invariantMass,
scattered_theta, scattered_p, proton_theta, proton_p,
positron_theta, positron_p, electron_theta, electron_p,
time_ee, x_FT, y_FT
\end{verbatim}
The data and both Monte Carlo arms must expose the same set, filled by the same code, or step~1
cannot be honoured.

\section*{Companion notes}

\begin{itemize}
\item \texttt{docs/RC\_METHOD\_FOR\_JPSI.md} --- how the $\pi^0$ RC-method study was done and what
to change for $J/\psi$.
\item \texttt{docs/RC\_METHOD\_COMPARISON.md} --- MC ratio versus the analytic \texttt{rc\_factor},
and what the difference costs.
\item \texttt{docs/rc\_method\_jpsi/rc\_method\_jpsi\_report.tex} --- the measured $\crad$ for this
campaign, both methods.
\item \texttt{docs/BORN\_XSEC\_EXTRACTION\_JPSI.md} --- the Markdown form of this note.
\end{itemize}

\end{document}
